\subsection*{Actual normalizer recognition on the genus-two target at nonsplit primes}

\paragraph{Theorem.}
Let $X_5:Y^2=1-X^5$ be its smooth proper genus-two model over a
number field containing $\mathbf Q(\zeta_5)$, let $O=\infty$ be its
unique point at infinity, and put $U_5=X_5\setminus\{O\}$.
For every prime $\ell=5$ or $\ell\not\equiv1\pmod5$, a nontrivial
closed procyclic subgroup of
\[
 P=\pi_1(U_{5,\overline K})^{(\ell)}
\]
with open arithmetic normalizer image is contained in an actual
cusp inertia subgroup. The assertion remains valid over every
finite number-field extension. The split primes
$\ell\equiv1\pmod5$ and the full profinite A.22 requirement are
not settled by this theorem.

\paragraph{1. The exact absence of a rational line at these primes.}
By F.3.19 the Jacobian $B=J_{X_5}$ is absolutely simple and has
an embedding $E=\mathbf Q(\zeta_5)\subset\operatorname{End}^0_K(B)$.
The full geometric endomorphism algebra is a division algebra
$D$, and $H_1(B(\mathbf C),\mathbf Q)$ is a faithful finite-dimensional
left $D$-vector space. Thus $[D:\mathbf Q]$ divides4. Since it
already contains the degree-four field $E$, one has
\[
 \operatorname{End}^0_{\overline K}(B)
 =\operatorname{End}^0_L(B)=E
\]
for every finite extension $L/K$.
Faltings' semisimplicity and endomorphism theorem, original
\S5,S\"atze3--4,p.360, identify
\[
 \operatorname{End}_{G_L}(V_\ell B)=E\otimes\mathbf Q_\ell.
\]
The module $V_\ell B$ is free of rank one over this algebra, by
the endomorphism-compatible Betti--etale comparison and the
degree-four action on four-dimensional rational homology.

For $\ell\ne5$, each factor of $E\otimes\mathbf Q_\ell$ has degree
$f=\operatorname{ord}_5(\ell)$: the primitive fifth roots have
Frobenius orbits of precisely that size over $\mathbf F_\ell$,
and the separable factors lift by Hensel's lemma. For $\ell=5$,
\[
 \Phi_5(T+1)=T^4+5T^3+10T^2+10T+5
\]
is Eisenstein, so the local algebra is a field of degree4.
At the primes in the theorem every field factor therefore has
degree greater than one. A nonzero idempotent on this rank-one
module has rank equal to a nonempty sum of those degrees, and
cannot have rank one. Semisimplicity would supply exactly such
an idempotent for any one-dimensional $G_L$-stable subspace.
There is consequently no rational Galois line even after an
arbitrary finite extension.

\paragraph{2. An actual finite cover detects the cyclic generator.}
Let
\[
 \Pi=\pi_1(U_5)/\ker(\pi_1(U_{5,\overline K})\to P),
 \quad p:\Pi\to G_K,\quad N=N_\Pi(A),\quad H=p(N).
\]
The open subgroup $H$ is also closed because $N$ is compact.
Since $P$ is a free pro-$\ell$ group of rank4, it is torsion-free
and $A\simeq\mathbf Z_\ell$. Choose a generator $u$ and a finite
$\ell$-group quotient $P\to Q$ in which its image $\bar u$ is
nonidentity. Put
\[
 P_0=\operatorname{preimage}(\langle\bar u\rangle)\subset P.
\]
The map $P_0\to\langle\bar u\rangle\to\mathbf Z/\ell$ takes
$u$ to a nonzero element. Its image in the finite-free
$\mathbf Z_\ell$-module $P_0^{\rm ab}$ is therefore nonzero.

The inverse image of $P_0$ in $\pi_1(U_{5,\overline K})$ is the
geometric covering subgroup of a connected finite etale curve
$Y_{\overline K}\to U_{5,\overline K}$. All equations and maps
descend to some finite extension $L/K$, and their finite-etale
property descends after enlarging $L$; the chosen geometric
connectedness is preserved. Its geometric maximal pro-$\ell$
group is exactly $P_0$. Indeed the Galois closure of this cover
has an $\ell$-group as deck group; the normal closure of a further
finite $\ell$-cover is obtained from the finitely many conjugate
$\ell$-covers over that Galois closure and again has an
$\ell$-group. This proves that every finite $\ell$-quotient of
$\pi_1(Y_{\overline K})$ factors through $P_0$, and proves the
maximal-pro-$\ell$ identification, not just a map onto it.

The arithmetic covering subgroup descends to an open subgroup
$\Pi_Y\subset\Pi$ with geometric kernel $P_0$ and projection
$G_L$. It does so because the kernel already divided out in
$\Pi$ is contained in its geometric covering subgroup. Set
$N_0=N\cap\Pi_Y$. Its index in $N$ is finite, hence
$p(N_0)$ is open and corresponds to a finite extension $L_0/L$.
The geometric kernel of $N_0\to p(N_0)$ lies in $P_0$ and acts
inner on $P_0$, so it acts trivially on $P_0^{\rm ab}$.
The nonzero homology line of $u$ is therefore an ordinary
$G_{L_0}$-stable line in the actual curve homology
\[
 P_0^{\rm ab}\otimes\mathbf Q_\ell
 =H_1(Y_{\overline K},\mathbf Q_\ell).
\]
No splitting of $N\to H$ or of $N_0\to G_{L_0}$ is used.

\paragraph{3. The same character has a first Magnus weight.}
The injective Magnus embedding of the rank-four free pro-$\ell$
group has a rational logarithm with a first nonzero homogeneous
term $u_d\in L_d(V_\ell B)$. This uses that $U_5$ has only one
geometric puncture, so its first homology is the proper homology
of $X_5$ and is pure of weight one. Every $L_d(V_\ell B)$ is
pure of weight $d$.

Conjugation by $n\in N$ acts as $u\mapsto u^{\theta(n)}$.
Geometric conjugation acts trivially on the first nonzero
homogeneous term, so comparison there shows that $\theta=1$
on $N\cap P$. Thus $\theta$ factors through $H$, and
\[
 \rho(p(n))u_d=\theta(n)u_d.
\]
This is the same character as on the detected homology line
of Step2. The assertion is an equality of actions on a
nonzero vector, with a torsion-free coefficient field, so
there is no finite-character ambiguity at this stage.

\paragraph{4. Weight two forces the actual boundary character.}
Let $\overline Y$ be the smooth proper compactification and
$T=\overline Y\setminus Y$. The homology localization sequence is
\[
 0\longrightarrow W_T\longrightarrow
 H_1(Y_{\overline K},\mathbf Q_\ell)
 \longrightarrow H_1(\overline Y_{\overline K},\mathbf Q_\ell)
 \longrightarrow0,
\]
where
\[
 W_T\simeq\left(\bigoplus_{t\in T(\overline K)}
                    \mathbf Q_\ell(1)e_t\right)
          \big/\mathbf Q_\ell(1)\Bigl(\sum_t e_t\Bigr).
\]
Take a good arithmetic Frobenius for both curves outside the
finite set of excluded places. A nonzero image of the detected
line in proper homology would have eigenvalue absolute value
$q^{1/2}$. Comparison with its first Magnus term, whose
absolute value is $q^{d/2}$, would force $d=1$. But then
$u_1$ would be a one-dimensional Galois-stable subspace of
$V_\ell B$ on an open subgroup, excluded by Step1.
Therefore the detected line has zero image in proper homology
and lies in $W_T$.

Choose a finite extension $L_1/L_0$ over which every point of
$T$ is rational. The displayed permutation representation
then becomes a direct sum and quotient of copies of the
same character $\chi_\ell$. Thus $G_{L_1}$ acts on all of
$W_T$ exactly by $\chi_\ell$, including on the nonzero
detected line. This proves the actual equality
$\theta=\chi_\ell$ on the inverse image of $G_{L_1}$ in $N_0$.
It also shows that the first Magnus degree was2. No
semisimple splitting of the mixed homology extension is
needed: the line was first placed in its actual kernel.

\paragraph{5. Passage back to the full arithmetic normalizer.}
The inverse image $N_1\subset N_0$ of $G_{L_1}$ surjects
onto $G_{L_1}$. Every element of the full normalizer in the
arithmetic group over $L_1$, namely
$N\cap p^{-1}(G_{L_1})$, has the same arithmetic image as
some element of $N_1$. Their quotient lies in $N\cap P$,
where Step3 proved that $\theta$ is trivial. Consequently
the actual full normalizer over $L_1$ acts on $A$ by
$\chi_\ell$ and has open (indeed surjective) arithmetic image.
Nakamura1994,Theorem2.1.1,p.94, for the full class of finite
$\ell$-groups now gives an actual cusp inertia subgroup
containing $A$. The finite field extensions leave the
geometric fundamental group and its unique-cusp conjugacy
class unchanged. This proves the theorem.

\paragraph{Exact remaining scope.}
At a split prime $\ell\equiv1\pmod5$, the four rank-one
CM eigenspaces are actual Galois lines, so Step1 supplies
no exclusion of the degree-one case. The theorem does not
assert that one of these lines lifts to a procyclic
arithmetic-normalizer subgroup. It specifies exactly the
prime range in which that possibility has been excluded.
